\chapter{图像复原、增强与压缩}
\label{chap:vision-image-restoration}

夜间拍到的招牌偏暗，字画周围有噪点，放大后又出现模糊和方块。把照片处理得明亮、锐利，并不能单凭外观判断招牌上的字已经恢复正确；把处理结果保存成较小文件，也不能只看重建图像就知道实际节省了多少存储。本章讨论受损图像怎样变得可用，并区分内容估计、呈现调整与表示成本三个目的。

设希望恢复的图像为$\bm{X}$，实际观测为$\bm{Y}$。一个便于分析的成像模型是
\begin{equation}
 \bm{Y}=\mathcal{Q}\bigl(\mathcal{S}(\bm{K}*\bm{X})+\bm{N}\bigr)\text{，}
 \label{eq:vision-image-degradation}
\end{equation}
其中$\bm{K}$描述模糊，$*$表示卷积，$\mathcal{S}$表示采样，$\bm{N}$表示噪声，$\mathcal{Q}$概括量化等后处理。它是一个简化模型：传感器噪声可能发生在颜色形成之前，图像也可能经历多次缩放和压缩，不能把这些操作一律交换顺序。RGB图像通常已经经过颜色插值、白平衡和色调映射；RAW保留更接近传感器读出的数据，两者可供反演的信息不同。

\term{图像复原}（image restoration）依据受损观测估计目标内容；\term{图像增强}（image enhancement）调整亮度、对比或颜色，使呈现适合用途；\term{图像压缩}（image compression）用更少比特表达图像。增强可以接在复原之后，压缩也可以与两者联合设计，但各自允许改变的量需要先规定。例如阅读招牌时允许提亮，档案记录却还需要保留原始观测和处理关系。下文的像素、信号和码流数字均为教学构造。

\section{图像复原}
\label{sec:vision-restoration-restoration}

恢复的困难取决于观测丢失了什么。噪声扰动像素，模糊混合邻域，下采样合并多个位置，缺口则完全没有相应观测。已知退化时，可以显式比较候选图像经过退化后是否接近输入；\term{盲复原}（blind restoration）是在使用时不知道全部退化参数，需要从图像及训练先验中共同估计。这里的“盲”描述可用信息，不能解释为已经覆盖所有真实相机和损坏形式。

\subsection{噪声干扰下的复原}
\label{sec:vision-restoration-denoising}

细小字画与噪声都可能表现为局部亮暗变化。要学会保留前者、去掉后者，训练需要说明噪声怎样产生，以及哪幅图像可作为清晰参考。最简单的模型是$\bm{Y}=\bm{X}+\bm{N}$，各元素$n_{i,j,k}$独立服从均值为零、方差为$\sigma^2$的高斯分布；真实传感器噪声可能随亮度变化，经过成像管线后也可能具有空间和通道相关性。

\subsubsection{DnCNN的噪声残差预测}
\label{sec:vision-restoration-dncnn}

DnCNN让网络预测观测与参考的差，而不是直接输出清晰图像。记这一残差网络为$r_{\bm{\theta}}$，由$n$对带噪图像与清晰参考训练
\begin{equation}
 L(\bm{\theta})=\frac{1}{2n}\sum_{i=1}^{n}
 \bigl\lVert r_{\bm{\theta}}(\bm{Y}_i)
             -(\bm{Y}_i-\bm{X}_i)\bigr\rVert_{\mathrm{F}}^2,
 \qquad
 \widehat{\bm{X}}=\bm{Y}-r_{\bm{\theta}}(\bm{Y})\text{。}
 \label{eq:vision-dncnn}
\end{equation}
这里的Frobenius范数表示所有像素与通道的平方和。网络以同尺寸卷积形成局部特征，中间层组合卷积、批归一化与ReLU，末层回到输入通道数，输出可正可负的残差。零填充维持空间大小，使相减能逐位置执行；它同时规定了图像边缘的处理方式\scite{vision-dncnn2017}

以招牌一行的四个灰度值为例，参考为$(0.2,0.2,0.8,0.8)$，加噪后得到$(0.25,0.17,0.84,0.74)$，残差标签就是$(0.05,-0.03,0.04,-0.06)$。若网络预测$(0.04,-0.02,0.03,-0.05)$，相减得到$(0.21,0.19,0.81,0.79)$。若网络把明暗跳变也解释成噪声，输出会抹掉字画；残差形式本身没有提供“噪声”和“真实细节”的语义区分，这一区分来自数据及损失。

固定噪声版本围绕给定$\sigma$合成配对，盲高斯版本在规定范围内抽取多个噪声强度。训练更新卷积参数与归一化统计，使用时保留这些参数和预处理，不再需要清晰参考。多任务的DnCNN-3还混合去噪、超分辨率和JPEG去块效应，但每类残差来自不同退化。把所有任务写成$\bm{Y}-\bm{X}$只是统一输出接口，不能推得真实噪声与JPEG伪影具有同一分布。

平方误差也决定了输出倾向。对固定观测$\bm{Y}=\bm{y}$和任意候选$\bm{a}$，有
\begin{equation}
 \mathbb{E}\!\left[\lVert\bm{X}-\bm{a}\rVert_{\mathrm{F}}^2
                  \mid\bm{Y}=\bm{y}\right]
 =
 \mathbb{E}\!\left[\lVert\bm{X}-\bm{\mu}\rVert_{\mathrm{F}}^2
                  \mid\bm{Y}=\bm{y}\right]
 +\lVert\bm{a}-\bm{\mu}\rVert_{\mathrm{F}}^2\text{，}
 \label{eq:vision-conditional-mean}
\end{equation}
其中$\bm{\mu}=\mathbb{E}[\bm{X}\mid\bm{Y}=\bm{y}]$，并假设二阶矩有限。展开平方后，交叉项因条件均值为零而消失，所以理想平方误差预测是条件均值。若一个细节存在多个可能位置，平均结果可能较平滑；这为后文比较像素保真与生成细节提供了依据。

\subsection{模糊观测下的复原}
\label{sec:vision-restoration-deblurring}

模糊不是逐像素独立扰动。相机运动或失焦使一个位置的光贡献到多个位置，恢复时要拆开这些混合。已知模糊核提供明确的反演关系；核未知时，学习式映射需要从数据识别可能的清晰结构。两条路线都依赖边界、对齐及退化范围。

\subsubsection{已知模糊核的频域正则化反卷积}
\label{sec:vision-restoration-deconvolution}

先考虑一维信号，假设模糊空间不变，并采用周期边界。卷积可写为$\bm{y}=\bm{H}\bm{x}+\bm{n}$，其中$\bm{H}$是循环矩阵。离散傅里叶变换将它对角化，每个频率满足$Y_\omega=H_\omega X_\omega+N_\omega$。直接除以$H_\omega$时，噪声变为$N_\omega/H_\omega$；弱频率响应会被大幅放大，零响应处则没有唯一反演。

为限制这种放大，可以把观测一致性与相邻位置的平滑要求合在一起：
\begin{equation}
 \widehat{\bm{x}}=\argmin_{\bm{x}}
 \lVert\bm{H}\bm{x}-\bm{y}\rVert_2^2+
 \lambda\lVert\bm{D}\bm{x}\rVert_2^2\text{，}
 \qquad \lambda>0\text{，}
 \label{eq:vision-deconv-objective}
\end{equation}
其中$\bm{D}$计算周期相邻差分。令梯度为零得到
$(\bm{H}^{\top}\bm{H}+\lambda\bm{D}^{\top}\bm{D})\widehat{\bm{x}}=\bm{H}^{\top}\bm{y}$。在二者均可由同一傅里叶基对角化且分母非零时，每个频率独立求解为
\begin{equation}
 \widehat X_\omega=
 \frac{\overline{H_\omega}}
 {|H_\omega|^2+\lambda|D_\omega|^2}\,Y_\omega\text{。}
 \label{eq:vision-regularized-inverse}
\end{equation}
上划线表示复共轭。惩罚较大的高频受到更强抑制，因而输出既减少噪声放大，也可能损失尖锐边缘。若$H_\omega=D_\omega=0$，此目标在该频率仍没有唯一解，需要另外规定约束。

取八点边缘$\bm{x}=(0,0,0,0,1,1,1,1)$，模糊核在中心和左右邻点的权重为$(0.6,0.2,0.2)$，再加入交替的$\pm0.04$噪声，得到观测
\[
 \bm{y}=(0.24,-0.04,0.04,0.16,0.84,0.96,1.04,0.76)\text{。}
\]
最高频处$H_\omega=0.2$，直接求逆把噪声幅值放大到0.2；用相邻差分与$\lambda=0.05$时，该处的逆增益为$0.2/(0.04+0.2)\approx0.833$，噪声分量幅值降为约0.033。

图\ref{fig:vision-deconvolution}显示全部八点结果。正则化后的过冲变小，但跃迁也变缓，且仍有轻微振铃。这里保留了负值与大于1的值以展示反演行为，没有在绘图前裁剪。真实图像若采用反射边界，不能继续照搬循环卷积公式；若模糊核错误或随位置变化，正则化只能缓解数值不稳定，不能修正错误的成像模型。

\begin{figure}[htbp]
 \centering
 % Eight-point circular blur, kernel (.2,.6,.2), alternating noise .04.
% DFT ridge inverse uses first differences and lambda=.05; values are un-clipped.
\begin{tikzpicture}[x=11mm,y=30mm,font=\figurefont\small]
 \draw[BookMuted,line width=.7pt,->] (-.35,-.3)--(7.5,-.3)
   node[right] {$i$};
 \draw[BookMuted,line width=.7pt,->] (-.35,-.3)--(-.35,1.38)
   node[above] {灰度};
 \foreach \y in {0,0.5,1} {
  \draw[BookRule,line width=.5pt] (-.35,\y)--(7.2,\y);
  \node[anchor=east] at (-.45,\y) {$\y$};
 }
 \foreach \i in {0,...,7} {
  \draw[BookMuted,line width=.5pt] (\i,-.3)--(\i,-.34);
  \node[below] at (\i,-.34) {$\i$};
 }
 \draw[BookInk,line width=1.2pt]
  plot coordinates {(0,0)(1,0)(2,0)(3,0)(4,1)(5,1)(6,1)(7,1)};
 \draw[BookAmber,dashed,line width=1pt]
  plot[mark=square*,mark size=1.5pt] coordinates
  {(0,.2)(1,-.2)(2,.2)(3,-.2)(4,1.2)(5,.8)(6,1.2)(7,.8)};
 \draw[BookTeal,line width=1pt]
  plot[mark=*,mark size=1.5pt] coordinates
  {(0,.1997)(1,-.0766)(2,-.01)(3,.133)(4,.867)(5,1.01)(6,1.0766)(7,.8003)};
 \draw[BookInk,line width=1.2pt] (0,1.65)--(.6,1.65);
 \node[anchor=west] at (.7,1.65) {参考};
 \draw[BookAmber,dashed,line width=1pt] (2.1,1.65)--(2.7,1.65);
 \node[anchor=west] at (2.8,1.65) {直接求逆};
 \draw[BookTeal,line width=1pt] (4.8,1.65)--(5.4,1.65);
 \node[anchor=west] at (5.5,1.65) {正则化};
\end{tikzpicture}

 \caption{八点周期边缘的教学反卷积。实线为清晰参考，带方块的虚线为直接求逆，带圆点的实线为差分正则化结果。两端因周期边界相接；纵轴未裁剪，便于观察过冲与平滑的取舍。}
 \label{fig:vision-deconvolution}
\end{figure}
\FloatBarrier

\subsubsection{Restormer的局部上下文与通道注意力复原}
\label{sec:vision-restoration-restormer}

运动模糊招牌的字画可能需要较大范围的结构信息。Restormer以多尺度编码器—解码器组织特征，把注意力矩阵放在通道维度上，控制高分辨率下的计算量。推理时不需输入显式模糊核，约束主要来自退化图像与清晰参考的训练配对\scite{vision-restormer2022}

输入经卷积得到特征，编码阶段逐级降低空间分辨率并增加通道，解码阶段上采样并拼接相应编码特征。在一个注意力头内，设位置数$n=hw$，通道数$d$。对归一化特征先作$1\times1$卷积混合通道，再作$3\times3$深度卷积混合邻域，得到按“通道$\times$位置”排列的$\bm{Q},\bm{K},\bm{V}\in\mathbb{R}^{d\times n}$。沿空间维归一化查询与键后，一种与实现一致的行向量记法是
\begin{equation}
 \bm{A}=\operatorname{softmax}_{\mathrm{row}}
       (\beta\,\bar{\bm{Q}}\bar{\bm{K}}^\top),
 \qquad
 \bm{O}=\bm{A}\bm{V}\text{，}
 \label{eq:vision-restormer-attention}
\end{equation}
其中$\beta$是每头学习的缩放量。$\bm{A}\in\mathbb{R}^{d\times d}$比较的是通道在整幅特征图上的响应关系，随后在每个位置混合通道；它不是每个像素直接对所有其他像素作加权求和。局部深度卷积及多尺度计算负责传播空间信息，全图统计又使通道组合依赖当前输入。

若$n=256\times256$而$d=32$，逐位置注意力需为一个头建立约$4.3\times10^9$个关系，通道矩阵则含1024项。仅考虑上述矩阵乘法，通道注意力的时间量级为$O(nd^2)$，关系矩阵存储为$O(d^2)$；还需要$O(nd)$的特征存储，以及卷积和其余网络的计算。说它对位置数近似线性，必须固定通道和其他结构。

门控前馈模块把归一化特征送入两个含逐点与深度卷积的支路，对一侧施加GELU后逐元素相乘，再投影并残差相加。乘法让一个支路调制另一个支路的响应，但门值不是必须位于$[0,1]$的概率。经过高分辨率细化后，末层预测图像修正量，与输入相加。这个残差符号与式\eqref{eq:vision-dncnn}中预测噪声再相减的约定不同。

训练以配对参考的像素损失优化网络，并逐渐增加裁剪块大小、相应减小批量，使模型在训练后期读取更大范围。对于招牌，小块可约束笔画，较大块还包含字间结构和周围边缘；使用时仍需核查笔画是否真实存在。去噪、去雨、运动模糊与失焦使用各自的数据和模型设置，同一个架构名不能代替检查点的适用范围。

\subsection{分辨率不足下的复原}
\label{sec:vision-restoration-super-resolution}

\term{超分辨率}（super-resolution）根据较低分辨率的输入估计更高分辨率的图像。放大网格增加了输出位置，却没有自动增加独立观测。例如把两个像素平均成一个时，$(0,1)$与$(0.5,0.5)$都会得到0.5。输出中的细线和文字怎样出现，取决于参考监督、退化假设与图像先验，而非输出像素数本身。

\subsubsection{插值输入与残差监督：DnCNN的超分辨率变体}
\label{sec:vision-restoration-dncnn-sr}

复用残差网络前，需要先让输入与目标处在同一网格。由高分辨率参考$\bm{X}$合成低分辨率图像$\bm{Y}=\mathcal{S}(\bm{X})$，再用双三次插值$\mathcal{U}$放大为$\bm{B}=\mathcal{U}(\bm{Y})$。DnCNN的相应变体学习$\bm{B}-\bm{X}$，使用时得到$\widehat{\bm{X}}=\bm{B}-r_{\bm{\theta}}(\bm{B})$。残差现在含插值和平滑带来的结构差异，不再是独立高斯噪声。

为了看清监督的含义，暂用两点平均和常数插值作教学简化。原细线为$(0,1)$，观测为0.5，插值得到$(0.5,0.5)$，残差标签为$(0.5,-0.5)$。若训练先验判断细线在右侧，预测残差能恢复这一排列；若真实原图为$(1,0)$，相同观测需要相反残差。唯一低分辨率观测无法区分两者。

像素监督衡量输出与训练参考的差。若应用还要求输出服从特定退化过程，应计算$\mathcal{S}(\widehat{\bm{X}})$并与$\bm{Y}$比较，这称为对观测的回投影核验。上述两个排列都能通过平均采样的回投影，所以观测一致性是必要依据，却不足以确认缺失细节。比较真实结果时还应固定放大倍数、采样滤波器和边界裁剪。

\Needspace{60mm}
\subsubsection{Real-ESRGAN的高阶退化合成与细节学习}
\label{sec:vision-restoration-real-esrgan}

网上照片可能先缩小、保存成JPEG，再截取、放大和重新压缩。只用一次双三次下采样训练，难以覆盖这种输入。Real-ESRGAN从清晰图像反复采样模糊、缩放、噪声和JPEG压缩，并使用sinc滤波模拟锐利边缘附近的振铃；原方法采用两轮退化构成训练配对\scite{vision-real-esrgan2021}

可以将一对训练样本写成$\bm{Y}=\mathcal{D}_2(\mathcal{D}_1(\bm{X}))$与$\bm{X}$。每轮内部都有随机参数，缩放可以向上或向下，噪声和压缩强度也变化，末尾操作另有顺序安排。高阶指重复退化过程，不表示对图像求高阶导数。这样得到的低清招牌仍与同一张清晰参考对应，监督目标可直接计算。

生成器沿用ESRGAN的残差密集结构。先以像素距离训练复原网络，再加入感知损失和对抗损失；前者比较固定视觉网络中的特征，后者让生成器面对区分真实和生成内容的判别器。U形判别器在多个局部位置给出反馈，并采用谱归一化约束训练。训练要更新生成器和判别器，使用时只保留生成器及相同输入约定，一次前向得到放大图像。

用多次JPEG保存的招牌检查结果时，应分别看方块是否减弱、笔画是否连贯、文字是否改变。重复退化扩大了合成覆盖，感知和对抗目标又鼓励自然细节；这些条件允许模型补出看似合理的纹理。它们既没有穷尽所有相机退化，也没有为补出的字画提供新的观测。

\subsubsection{DiffBIR的退化移除与生成先验补充}
\label{sec:vision-restoration-diffbir}

严重退化下，直接复原容易产生较平滑的结果。DiffBIR将处理拆为减轻退化与补充内容两个阶段：前阶段提供较稳定的图像条件，后阶段用预训练扩散模型补充细节。本节采用2024年3月的第二版预印本描述，区分其中的训练条件构造和实际使用\scite{vision-diffbir2024}

训练生成模块时，作者另训练一个像素平方误差复原器，在较宽的简化退化范围内产生条件图像$\bm{X}_{\mathrm{c}}$。这个预处理器的职责是提供多样、较少退化干扰的训练条件；实际使用时改接各任务已有的复原模块，例如面部复原或盲超分辨率模块。把这两者统称为“第一阶段”会掩盖一个重要差别：训练时的条件产生器不必是交付时处理真实照片的复原器。

生成阶段用固定的图像编码器把$\bm{X}_{\mathrm{c}}$变为条件潜变量$\bm{z}_{\mathrm{c}}$。IRControlNet复制预训练去噪网络的编码及中间模块，读取$\bm{z}_{\mathrm{c}}$与当前带噪潜变量的通道拼接。支路经零初始化卷积输出修正，加入固定去噪网络的中间与跳连特征。训练给清晰参考潜变量加噪，以噪声预测平方误差只更新条件网络及连接参数；主去噪网络、图像编码器和解码器保留预训练参数。共同的潜空间条件生成接口将在\ref{sec:vision-models-media-generation}节集中说明。

推理从随机潜变量开始，反复用条件支路与固定生成先验去噪，再解码成图像。可选的复原引导在每步得到清晰潜变量估计$\tilde{\bm{z}}^{(t)}$后，减少解码图像与前阶段结果的加权距离：
\begin{equation}
 \begin{aligned}
 L_{\mathrm{g}}(\tilde{\bm{z}})
 &=\frac{1}{hwc}
   \left\lVert\bm{W}\odot
    \bigl(\mathcal{D}(\tilde{\bm{z}})-\bm{X}_{\mathrm{c}}\bigr)
   \right\rVert_{\mathrm{F}}^2,\\
 \tilde{\bm{z}}_{\mathrm{g}}^{(t)}
 &=\tilde{\bm{z}}^{(t)}
   -s\nabla_{\tilde{\bm{z}}}L_{\mathrm{g}}
       (\tilde{\bm{z}}^{(t)})\text{。}
 \end{aligned}
 \label{eq:vision-diffbir-guidance}
\end{equation}
这里$\mathcal{D}$是图像解码器，$s$控制引导强度，$\bm{W}$根据前阶段图像的局部梯度统计形成，对较平坦区域赋予更大权重。更新后的清晰估计参与下一采样步。损失中的权重在平方之内，实际像素平方误差的系数是$w_{i,j,k}^2$。

在固定低清人像和随机种子下，增加引导会加强对前阶段轮廓及平坦区域的约束，同时可能压低生成细节。这里被追随的仍是一个复原估计，它若已把眼睛或文字恢复错，引导也可能保留错误。皮肤纹理、身份细节与招牌字形应分别核验，不能把新的高频外观当作相机曾经记录到的事实。

\subsection{缺失区域的复原}
\label{sec:vision-restoration-inpainting}

坏点或缺口与加性噪声不同：对应位置根本没有可信观测。以二值掩码$\bm{M}$标明有效像素，$m_{i,j}=1$表示可用，0表示缺失，并在通道维广播。恢复目标是依据周围内容估计缺口原有内容；若要求删除杯子并创造一只花瓶，则还包含新的语义目标，交由\ref{chap:vision-image-generation}章讨论。

\subsubsection{部分卷积的有效像素归一化与掩码更新}
\label{sec:vision-restoration-partial-convolution}

普通卷积会读取缺口中填入的零或均值，输出因占位数值而变化。\term{部分卷积}（partial convolution）先屏蔽无效位置，再按有效数量调整响应。把一个窗口的输入、掩码和权重分别展平为$\bm{x},\bm{m},\bm{w}\in\mathbb{R}^q$，其输出及新掩码为
\begin{equation}
 x'=
 \begin{cases}
 \dfrac{q}{\sum_i m_i}\bm{w}^{\top}(\bm{x}\odot\bm{m})+b,
   &\sum_i m_i>0,\\
 0,&\sum_i m_i=0,
 \end{cases}
 \qquad
 m'=\mathbb{I}\!\left[\sum_i m_i>0\right]\text{。}
 \label{eq:vision-partial-convolution}
\end{equation}
其中$b$是偏置。归一化因子只乘加权和，不乘偏置；完全无有效输入时输出零，新掩码也为零\scite{vision-partial-convolution2018}

取窗口$\bm{x}=(0.2,0,0.8)$、$\bm{m}=(1,0,1)$、$\bm{w}=(1/3,1/3,1/3)$且$b=0$，屏蔽后的加权和为$1/3$，再乘$3/2$得到0.5，$m'=1$。若把缺口占位改成100，结果仍为0.5。窗口内全缺失时则不能计算这种归一化，必须走零输出分支。图\ref{fig:vision-partial-mask}说明，堆叠窗口使可计算范围逐层进入缺口。

\begin{figure}[htbp]
 \centering
 % Binary masks; kernel width 3, stride 1. No image values are fabricated.
\begin{tikzpicture}[x=1mm,y=1mm,font=\figurefont\small,
 arr/.style={-{Stealth[length=2mm]},BookMuted,line width=.8pt}]
 \foreach \r/\textlabel in {0/输入掩码,1/第一层,2/第二层} {
  \node[anchor=east] at (0,{-\r*18}) {\textlabel};
  \foreach \i in {0,...,6} {
   \draw[BookMuted,line width=.6pt] ({5+\i*11},{-\r*18-4})
     rectangle ++(11,8);
  }
 }
 \foreach \r/\indices in {0/{0,1,5,6},1/{0,1,2,4,5,6},2/{0,1,2,3,4,5,6}} {
  \foreach \i in \indices {
   \fill[BookBlue!12] ({5+\i*11+.4},{-\r*18-3.6}) rectangle ++(10.2,7.2);
   \fill[BookBlue] ({10.5+\i*11},{-\r*18}) circle (.8mm);
  }
 }
 \draw[BookAmber,line width=1pt] (27,-5.5) rectangle (60,5.5);
 \draw[BookAmber,line width=1pt] (27,-23.5) rectangle (60,-12.5);
 \draw[arr] (43.5,-6)--(43.5,-12);
 \draw[arr] (43.5,-24)--(43.5,-31);
 \node[anchor=north] at (43.5,-43) {中心掩码：$0\to0\to1$};
\end{tikzpicture}

 \caption{宽度为3、步幅为1的掩码更新教学例。带点填充的格子表示已有有效输入，白格表示缺失；每层只要窗口内有一个有效位置，该层输出就标为有效。图中省略边界外输入。中间格在第二层获得可计算支持，不代表新取得了原像素观测。}
 \label{fig:vision-partial-mask}
\end{figure}

部分卷积修补网络把特征和掩码同时传过编码器、解码器及跳连。训练从清晰图像随机去掉不规则区域，分别约束已知区和缺失区的像素误差，并加入感知特征、风格统计及缺口周边平滑约束。感知项关心特征结构，风格项关心通道相关统计，它们都不能唯一确定遮挡背后的实际对象。

使用时输入受损图像和原始掩码，网络预测$\widehat{\bm{X}}_{\mathrm{net}}$，再形成
$\widehat{\bm{X}}=\bm{M}\odot\bm{Y}+(1-\bm{M})\odot\widehat{\bm{X}}_{\mathrm{net}}$，使已知像素按要求保留。宽度很小的划痕常有两侧结构可参考，大缺口则允许多种内容。中间掩码变为1只说明特征已有计算来源，原始观测掩码仍须保留，用来区分记录与估计。

\section{图像增强}
\label{sec:vision-restoration-enhancement}

有些图像已经保留了需要的结构，只是亮度或对比不适合观看；另一些暗图则同时缺少可靠观测。两种情况可能有相似外观，但只改变显示曲线与重新估计内容会带来不同结果。增强的目标应说明希望看清什么，以及哪些颜色、测量量或局部关系必须保留。

\subsection{暗处可见性与曝光的改善}
\label{sec:vision-restoration-low-light}

低照度时，传感器收到的光子较少，信号相对噪声较弱。若将像素与加性噪声一起乘上$a$，信号幅值变为$a$倍，噪声标准差也变为$a$倍，这种线性提亮没有改善二者的比例。延长真实曝光能收集更多光子，却可能引入运动模糊；对已经拍下的短曝光照片，事后乘增益不能替代重新曝光。

\subsubsection{Learning to See in the Dark的RAW到RGB映射}
\label{sec:vision-restoration-sid}

Learning to See in the Dark直接读取短曝光RAW，并用同场景长曝光图像形成的RGB作为监督。它把颜色采样整理、去噪和颜色形成放在同一学习流程中，因而可利用RGB图像形成之前的信息\scite{vision-sid2018}

以Bayer传感器为例，每个$2\times2$位置采样一红、两绿和一蓝，原始阵列为$h\times w\times1$。按位置打包后得到$(h/2)\times(w/2)\times4$数组；两个绿色位置保留为不同通道。打包只是重新排列观测，没有提前补齐每个像素的三种颜色。输入减去传感器黑电平，按有效范围归一化，再乘外部给定的放大比例。网络用编码器—解码器预测半分辨率的12通道结果，经子像素重排形成全分辨率三通道RGB。

假设某个RAW读数为530，黑电平为512，可用上界为16383，则归一化信号为$(530-512)/(16383-512)\approx0.00113$。乘100后约为0.113。若忽略黑电平，直接用$530/16383$再乘100会超过3，数值主要来自偏置并很快被裁剪；这不是同一种输入。相机的黑电平、颜色阵列和增益约定必须与训练一致。

训练以短曝光输入与对齐长曝光RGB参考之间的$L_1$距离更新网络。曝光比为放大比例提供依据，使用时放大比例也可作为外部控制，原方法并不自动推断最佳曝光。长曝光图像仍可能有噪声，移动人物和树叶也会破坏配对；若参考错位，逐像素监督就会把位置差异也当作待修正误差。

固定同一夜景比较常规RGB、RGB提亮与RAW映射时，应查看暗处结构、颜色和亮部饱和，并记录三者实际使用的数据。RAW路线拥有更早阶段的传感器信息，不能把收益全部归结为更强的RGB后处理。相机更换和动态场景还需要单独验证，静态配对训练没有自动解决运动。

\subsection{色彩、对比与层次的调整}
\label{sec:vision-restoration-tone}

当图像已有可用内容时，可以直接规定数值映射，不必训练网络。全局调整让相同输入值得到相同输出，局部调整则依据邻域分布改变映射；后者更灵活，也更容易改变区域之间的相对关系。

\subsubsection{白平衡与全局色调曲线}
\label{sec:vision-restoration-white-balance}

\term{白平衡}（white balance）按颜色通道调整增益，使选定照明条件下的中性物体呈现为中性色。在简化的线性RGB模型中，$\bm{c}'=\operatorname{diag}(g_r,g_g,g_b)\bm{c}$。若灰色参考被拍成$(0.20,0.16,0.10)$，取增益$(0.8,1,1.6)$会得到$(0.16,0.16,0.16)$。这依赖参考本来应中性；把一面黄色墙误当灰卡会使整图颜色偏移。色彩管理还包含颜色空间变换等环节，通道增益只表达其中的白平衡近似。

色调曲线把亮度$I\in[0,1]$映射成$T(I)$。为避免不同软件对“伽马”的方向约定混淆，这里直接规定$T(I)=I^\gamma$。当$\gamma=1/2$时，0.04、0.25和0.81依次变为0.2、0.5和0.9，暗部变化较大；$\gamma>1$则压暗中间值。若先作两倍线性提亮再裁剪，0.6和0.9都变为1，高光之间的差异就丢失了。

逐通道应用非线性曲线通常会改变通道比例。若希望尽量保持线性RGB比例，可先计算亮度$I$，再在$I>0$处用$\bm{c}'=\bm{c}\,T(I)/I$统一缩放，随后处理超出色域的数值；接近零时需稳定分母并规定黑色输出。保持比例不等于所有显示条件下的感知色相严格不变，颜色空间和裁剪仍有作用。

这些操作重新分配现有层次。招牌暗部可能更容易辨认，平坦暗区的噪声也可能一起显露；已经被传感器饱和或量化合并的字画，不能仅靠单调曲线恢复。

\subsubsection{受限局部直方图均衡化}
\label{sec:vision-restoration-clahe}

同一全局曲线不容易同时照顾阴影和高光。\term{受限局部直方图均衡化}（contrast limited adaptive histogram equalization, CLAHE）为图像分块建立亮度直方图，限制过高的计数峰，并用重新分配后的累计分布构造局部映射。限制峰值是为了避免把少量灰度差异拉伸得过强；在相邻块的映射结果之间插值，则减少块边界突变。

用四个亮度箱说明计算。某块的直方图为$(0,8,0,0)$，共有8个像素，若直接均衡化，第二箱的累计概率立刻达到1，微小变化会引起很大的映射变化。令计数上限为2，截断后得到$(0,2,0,0)$，需要重新分配6个计数。以保持上限为目标，将这些计数均匀补入其余未满的箱，得到$(2,2,2,2)$，累计概率为$(0.25,0.5,0.75,1)$。这个小例采用容量约束的整数分配；实现中的阈值归一化和余数分配规则须记录，不能只写一个无单位的“阈值2”。

一个像素通常受到周围四个块的映射影响。先将它的亮度分别输入四条曲线，再按位置双线性插值，因而即便亮度相同，不同位置的输出也可能不同。块较小时更关注局部，但噪声及小结构更容易主导直方图；上限放宽又增加局部对比。对招牌可以同时检查细字与平坦墙面，避免只看被增强的字画而忽略噪声和光晕。医学、档案和测量图像还应规定哪些灰度关系允许变化。

\section{图像压缩}
\label{sec:vision-restoration-compression}

压缩的输出是可传输的比特流。一个网络即使能把图像映射到很小的浮点数组，也不等于已经压缩：浮点数需要多少位、哪些辅助信息要传、解码器如何复现概率和符号，都必须明确。以下用尺度超先验方法连接表示学习与实际编解码。

\subsection{表示变换与量化}
\label{sec:vision-restoration-quantization}

分析变换$g_{\mathrm{a}}$把图像变为主潜变量$\bm{y}$，合成变换$g_{\mathrm{s}}$从离散潜变量恢复图像：
\begin{equation}
 \bm{y}=g_{\mathrm{a}}(\bm{X}),\qquad
 \hat y_i=\operatorname{round}(y_i),\qquad
 \widehat{\bm{X}}=g_{\mathrm{s}}(\hat{\bm{y}})\text{。}
 \label{eq:vision-transform-coding}
\end{equation}
\term{量化}（quantization）把连续数映射到有限精度的离散值，舍弃区间内的差别。此处以单位步长取整为例，步长的作用也可由分析与合成变换吸收。学习式变换通过卷积及非线性组织空间、通道信息，训练目标共同决定哪些变化值得保留。

用两个像素作可复算的变换例。对$\bm{x}=(0.2,0.8)$，令$y_1=2(x_1+x_2)=2$、$y_2=2(x_2-x_1)=1.2$。取整得到$(2,1)$，合成按$x_1=(y_1-y_2)/4$、$x_2=(y_1+y_2)/4$得到$(0.25,0.75)$。均方误差为0.0025。这里分析与合成是手工线性变换，学习式网络用数据寻找更适合图像的非线性变换；量化造成不可逆信息损失的原因相同。

取整函数的导数几乎处处为零，不能直接提供有用梯度。尺度超先验方法在训练中以独立均匀噪声替代取整，即$\tilde y_i=y_i+u_i$、$u_i\sim\mathcal{U}(-1/2,1/2)$。带波浪号的连续随机变量用于可微训练，带帽号的整数用于实际编码。训练近似不能代替交付时真正取整、编码并解码的检查。

设$R$为估计码长，$D$为重建失真，训练优化$R+\lambda D$。固定两者的归一化口径时，较大的$\lambda$让失真受到更重惩罚，网络可能保留更多细节并付出更多码率；这不是任意单幅图像上单调改善所有指标的保证。通常需要多个工作点形成码率—失真曲线，而非只调一个数就宣称适合全部用途。

\subsection{概率估计与码流形成}
\label{sec:vision-restoration-bitstream}

量化后还需给整数分配码字。常见符号应使用较短表示，少见符号可以较长，但编码器与解码器必须共享同一概率表。主潜变量在平坦区域常聚集于小幅值，在纹理区域可能更分散；一个全局分布难以充分表达这些变化。

\Needspace{45mm}
\subsubsection{尺度超先验的辅助表示与条件概率}
\label{sec:vision-restoration-hyperprior}

\term{尺度超先验}（scale hyperprior）为主潜变量的局部尺度建立辅助表示\scite{vision-hyperprior2018}
超分析变换产生$\bm{z}=h_{\mathrm{a}}(\bm{y})$，量化为$\hat{\bm{z}}$后，由超合成变换得到尺度$\bm{\sigma}=h_{\mathrm{s}}(\hat{\bm{z}})$。每个$\sigma_i>0$控制一个主符号的概率分布，辅助符号本身也需要编码。

原方法以条件零均值高斯描述主潜变量，并将密度在单位量化区间上积分。对整数$k$，离散概率为
\begin{equation}
 P_i(k\mid\hat{\bm{z}})
 =\int_{k-1/2}^{k+1/2}
   \frac{1}{\sqrt{2\pi}\sigma_i}
   \exp\!\left(-\frac{u^2}{2\sigma_i^2}\right)\,\mathrm{d}u
 =\Phi\!\left(\frac{k+1/2}{\sigma_i}\right)
  -\Phi\!\left(\frac{k-1/2}{\sigma_i}\right)\text{，}
 \label{eq:vision-hyperprior-probability}
\end{equation}
其中$\Phi$是标准正态累计分布函数。不能把密度在整数处的值直接当作符号概率；密度有单位，且在整数点求和不保证等于1。小$\sigma_i$给零附近符号较高概率，大$\sigma_i$为较大幅值留出更多概率。

辅助潜变量采用学习到的因子化分布，编码时也在量化区间上取得离散概率。一次图像的理想估计码长为
\begin{equation}
 R_{\mathrm{est}}=
 -\sum_j\log_2 P_j^{z}(\hat z_j)
 -\sum_i\log_2 P_i(\hat y_i\mid\hat{\bm{z}})\text{。}
 \label{eq:vision-hyperprior-rate}
\end{equation}
训练时相应使用加噪潜变量的连续密度，并与重建损失共同更新分析、合成及概率模型。辅助流的项不可省略，否则网络可以把大量信息转移到$\bm{z}$而不付出代价，目标就不再衡量总压缩成本。

假设共享全局分布编码主符号需要100 bit，传入尺度辅助信息后主流降至72 bit。若辅助流为12 bit，总计84 bit，有16 bit收益；若辅助流需要35 bit，总计107 bit，反而增大。只有节省的主流成本超过辅助成本时，这种条件化才划算。图\ref{fig:vision-hyperprior-codec}标出辅助信息如何同时让两端构造主流概率。

\begin{figure*}[htbp]
 \centering
 % Encoder/decoder dependencies. All ordinary edges are axis-aligned.
\begin{tikzpicture}[x=1mm,y=1mm,font=\figurefont\small,
 b/.style={draw=BookBlue,fill=BookBlue!8,minimum width=13mm,
   minimum height=8mm,inner sep=1mm,line width=.8pt,align=center},
 o/.style={b,draw=BookTeal,fill=BookTeal!8},
 s/.style={draw=BookMuted,minimum width=12mm,minimum height=8mm,
   inner sep=1mm,line width=.8pt},
 a/.style={-{Stealth[length=1.8mm]},BookInk,line width=.9pt},
 c/.style={a,BookAmber,rounded corners=1mm}]
 \node at (5,0) (x) {$\bm{X}$};
 \node[b] at (22,0) (ga) {$g_{\mathrm{a}}$};
 \node[b] at (42,0) (qy) {$Q_y$};
 \node[b] at (64,0) (ey) {E};
 \node[s] at (83,0) (sy) {$S_y$};
 \node[o] at (102,0) (dy) {D};
 \node[o] at (123,0) (gs) {$g_{\mathrm{s}}$};
 \node at (144,0) (xh) {$\widehat{\bm{X}}$};
 \node[b] at (22,-32) (ha) {$h_{\mathrm{a}}$};
 \node[b] at (42,-32) (qz) {$Q_z$};
 \node[b] at (64,-32) (ez) {E};
 \node[s] at (83,-32) (sz) {$S_z$};
 \node[o] at (102,-32) (dz) {D};
 \node at (123,-32) (zh) {$\hat{\bm{z}}$};
 \node[b] at (42,-16) (hse) {$h_{\mathrm{s}}$};
 \node[o] at (123,-16) (hsd) {$h_{\mathrm{s}}$};
 \foreach \u/\v in {x/ga,ga/qy,qy/ey,ey/sy,sy/dy,dy/gs,gs/xh,
   ha/qz,qz/ez,ez/sz,sz/dz,dz/zh} {\draw[a] (\u)--(\v);}
 \draw[a] (ga)--(ha);
 \draw[c] (qz)--(hse);
 \draw[c] (hse.east)--(64,-16)--(ey.south);
 \draw[c] (zh)--(hsd);
 \draw[c] (hsd.west)--(102,-16)--(dy.south);
 \node[anchor=west,text=BookAmber] at (66,-13) {$\bm{\sigma}$};
 \node[anchor=east,text=BookAmber] at (100,-13) {$\bm{\sigma}$};
 \node[anchor=south] at (33,4) {$\bm{y}$};
 \node[anchor=south] at (112.5,4) {$\hat{\bm{y}}$};
 \node[anchor=north,text=BookMuted] at (42,-40) {编码端};
 \node[anchor=north,text=BookMuted] at (112,-40) {解码端};
\end{tikzpicture}

 \caption{尺度超先验的实际编解码依赖。$Q$为取整，E与D分别为熵编码和熵解码，$S_y,S_z$为两条比特流。两端都从同一组整数$\hat{\bm{z}}$经$h_{\mathrm{s}}$产生尺度，主流解码因此必须等待辅助流。图中省略参数与尺寸元数据，未画入不可传输的训练噪声或清晰参考。}
 \label{fig:vision-hyperprior-codec}
\end{figure*}
\FloatBarrier

\subsubsection{量化符号的熵编码与实际码长}
\label{sec:vision-restoration-entropy-coding}

\term{熵编码}（entropy coding）根据共享概率把离散符号无损地转成比特。这里“无损”指整数符号能够原样恢复，前面的量化仍然有损。算术编码的一种解释是不断缩小$[0,1)$区间，用落在最终区间内的二进制前缀标识整个序列。

设三个教学符号$A,B,C$的概率依次为$1/2,1/4,1/4$，基础区间为$[0,1/2)$、$[1/2,3/4)$、$[3/4,1)$。编码序列$ABC$时，$A$把当前区间缩为$[0,1/2)$，$B$选择其中第二段得到$[1/4,3/8)$，$C$再选择最后四分之一得到$[11/32,12/32)$。其宽度为$1/32$，对应理想码长5 bit；二进制前缀$01011$恰好表示这个区间。

若解码端已知长度为3，可用区间内代表值$11/32$恢复序列。它落在基础区间的$A$段，归一化到该段后得到$11/16$，落在$B$段；再次减去$1/2$并除以$1/4$，得到$3/4$，落在$C$段。区间统一左闭右开，因此边界也有确定归属。图\ref{fig:vision-arithmetic-coding}画出逐步缩小的区间；实际实现使用有限精度整数和归一化输出位，无需保存无限精度实数。

\begin{figure}[htbp]
 \centering
 % All intervals use the same coordinate scale; probabilities .5,.25,.25.
\begin{tikzpicture}[x=80mm,y=1mm,font=\figurefont\small]
 \foreach \row/\lab in {0/初始,1/$A$,2/$AB$,3/$ABC$} {
  \draw[BookRule,line width=.7pt] (0,{-15*\row})--(1,{-15*\row});
  \node[anchor=east] at (-.05,{-15*\row}) {\lab};
 }
 \fill[BookBlue!15] (0,-3) rectangle (1,3);
 \draw[BookBlue,line width=.8pt] (0,-3) rectangle (1,3);
 \fill[BookTeal!20] (0,-18) rectangle (.5,-12);
 \draw[BookTeal,line width=.8pt] (0,-18) rectangle (.5,-12);
 \fill[BookTeal!35] (.25,-33) rectangle (.375,-27);
 \draw[BookTeal,line width=.8pt] (.25,-33) rectangle (.375,-27);
 \fill[BookTeal] (.34375,-48) rectangle (.375,-42);
 \foreach \x/\lab in {0/0,.5/{1/2},1/1} {
  \node[anchor=south] at (\x,4) {$\lab$};
 }
 \node[anchor=west] at (.54,-15) {$[0,1/2)$};
 \node[anchor=west] at (.43,-30) {$[1/4,3/8)$};
 \node[anchor=west] at (.43,-45) {$[11/32,12/32)$};
 \draw[BookMuted,dashed,line width=.6pt] (.5,-18)--(.375,-27);
 \draw[BookMuted,dashed,line width=.6pt] (0,-18)--(.25,-27);
 \draw[BookMuted,dashed,line width=.6pt] (.25,-33)--(.34375,-42);
 \draw[BookMuted,dashed,line width=.6pt] (.375,-33)--(.375,-42);
\end{tikzpicture}

 \caption{固定概率下序列$ABC$的算术编码教学例。每行共享$[0,1)$坐标，实色区间为读入当前前缀后仍可能的范围。最终区间宽$1/32$，可由5位前缀$01011$表示；序列长度和概率表在本例中预先共享，不计入这5位。}
 \label{fig:vision-arithmetic-coding}
\end{figure}

在尺度超先验模型中，要把式\eqref{eq:vision-hyperprior-probability}离散化成编码器可用的整数累计概率表，规定精度、尾部符号与逃逸机制，并使两端采用一致的符号顺序。零概率不能分配给实际出现的符号，否则无法编码。有限精度概率、流终止、填充以及尺寸和版本元数据都会使实际字节数偏离负对数概率估计。

真实文件应先提供足以解析的尺寸和模型约定，再存辅助流与主流，或用等价可解析的组织方式记录偏移。若一张$256\times256$图像的全部文件为4096 byte，则实际每像素比特数为$8\times4096/(256\times256)=0.5$，分母不再乘RGB的通道数。模型权重若由收发双方预先共享，常规每图码率不含权重；若每张图传一个专用模型，则这项成本也必须纳入所声明的传输预算。

\subsection{解码重建与用途保持}
\label{sec:vision-restoration-decoding}

解码按依赖执行。用共享的辅助概率模型从$S_z$恢复$\hat{\bm{z}}$，由$h_{\mathrm{s}}$产生主流概率，再从$S_y$恢复$\hat{\bm{y}}$，最后由$g_{\mathrm{s}}$形成$\widehat{\bm{X}}$。编码端也必须用量化后的$\hat{\bm{z}}$计算尺度，不能用未经传输的连续$\bm{z}$，否则两侧的概率表可能不同。

可以用一次独立解码检查这一点：只给解码程序文件、声明的模型参数和必要配置，确认它能复现输出。原图、训练参考和编码中间浮点值都不在它的可用输入中。一个只展示分析网络和合成网络前向结果的实验，尚未完成这个检查。

同样的0.5 bit/像素可以优先保留平滑颜色，也可以保留文本边缘，不同失真目标会形成不同取舍。应把解码图像送入固定的文字读取、识别或测量流程，检查招牌字符串、车牌边界及尺寸。感知上自然的纹理不一定有助于这些用途，较小的平均像素误差也可能掩盖一个关键数字的错误。视频中的运动流、残差流和参考帧依赖将在\ref{chap:vision-video-restoration}章接续。

\section{处理效果的验证与选择}
\label{sec:vision-restoration-evaluation}

判断处理是否值得，需要同时看内容、呈现与成本。比较前固定输入退化、参考对齐、像素范围和颜色约定；对没有可信参考的真实图像，应明确哪些结论只能来自任务结果或人工观察，不能给不存在的原图计算“恢复正确率”。

\subsection{内容保真与感知质量}
\label{sec:vision-restoration-quality}

参考指标逐对比较处理结果与对应图像，感知判断则关心人眼看到的结构或自然程度。两者可以相关，但比较对象不同，尤其在一幅低清观测对应多个可能细节时。

\subsubsection{PSNR与SSIM的像素和局部结构比较}
\label{sec:vision-restoration-psnr-ssim}

对含$n$个有效标量元素的对齐图像，均方误差和\term{峰值信噪比}（peak signal-to-noise ratio, PSNR）为
\begin{equation}
 \operatorname{MSE}=\frac1n\sum_{i=1}^{n}(x_i-\hat x_i)^2,
 \qquad
 \operatorname{PSNR}
 =10\log_{10}\frac{v_{\max}^2}{\operatorname{MSE}}\text{。}
 \label{eq:vision-psnr}
\end{equation}
$v_{\max}$是约定的峰值，归一化图像通常取1，8位整数通常取255；误差为零时PSNR为正无穷。前述两像素压缩例的MSE为0.0025，取$v_{\max}=1$得到约26.02 dB。对一组图像平均各图PSNR，与先合并所有像素误差再计算PSNR不是同一聚合方式。

\term{结构相似性}（structural similarity, SSIM）在对应局部窗口中比较亮度、对比和共同变化。其常用合并形式为
\begin{equation}
 \operatorname{SSIM}(\bm{x},\hat{\bm{x}})
 =
 \frac{(2\mu_x\mu_{\hat x}+c_1)(2\sigma_{x,\hat x}+c_2)}
 {(\mu_x^2+\mu_{\hat x}^2+c_1)
  (\sigma_x^2+\sigma_{\hat x}^2+c_2)}\text{，}
 \label{eq:vision-ssim}
\end{equation}
其中$\mu$是局部加权均值，$\sigma^2$是相应方差，$\sigma_{x,\hat x}$是协方差；$c_1,c_2>0$稳定近常量窗口的计算。整图值由局部结果聚合，窗口大小、权重、边界和颜色处理都应固定\scite{vision-ssim2004}

取参考$(0,0,1,1)$，平滑结果$(0,0.25,0.75,1)$的MSE为0.03125；只把第三个位置改成0的结果$(0,0,0,1)$，MSE为0.25，按固定长度比较时也可视为边缘平移一格。再取$(0,0,0.8,1)$，MSE只有0.01，但若第三位置恰是识别一笔关键字画的依据，任务错误仍可能很大。像素分数描述这些改变的平均幅度，SSIM还考察局部共同变化，二者都不知道这笔画的业务含义。

配准误差会使正确内容与错误位置比较。评价应先确认对应关系，再统一是否裁剪边缘、只评亮度还是评RGB，以及线性值或显示编码值。不能为每个方法单独选择最有利的对齐和颜色口径。

\Needspace{65mm}
\subsubsection{LPIPS的深层特征距离}
\label{sec:vision-restoration-lpips}

像素变化的大小未必与感知变化一致。\term{学习感知图像块相似度}（learned perceptual image patch similarity, LPIPS）使用视觉网络的多层特征，并依据人类对图像块相似性的判断校准组合权重\scite{vision-lpips2018}

设第$\ell$层在位置$p$的特征向量经通道归一化后为$\bar{\bm{f}}_{\ell,p}(\bm{X})$，该层共有$n_\ell$个位置，学习到的通道权重为$\bm{w}_\ell$。其距离写成
\begin{equation}
 d_{\mathrm{LPIPS}}(\bm{X},\widehat{\bm{X}})
 =
 \sum_\ell\frac1{n_\ell}\sum_p
 \left\lVert
 \bm{w}_\ell\odot
 \bigl(\bar{\bm{f}}_{\ell,p}(\bm{X})
      -\bar{\bm{f}}_{\ell,p}(\widehat{\bm{X}})\bigr)
 \right\rVert_2^2\text{。}
 \label{eq:vision-lpips}
\end{equation}
两幅图像都经过同一网络与预处理。先按位置比较归一化特征，再平均空间差异并累加层，不能将整图分类向量的余弦距离直接叫作LPIPS。若某位置的两通道差为$(0.2,0.1)$，权重为$(1,2)$，该位置贡献为$0.2^2+(2\times0.1)^2=0.08$。实际值取决于骨干与校准版本。

平滑砖墙可能产生较小像素误差，却改变深层纹理特征；把招牌“8”改成“3”又可能只影响很小区域，整体感知距离仍然较小。LPIPS衡量对应图像的特征差异，不能独立核验文字、身份或事实，也不能把一项分数直接解释成“恢复了多少真实内容”。

感知与失真的理论取舍采用另一个更具体的定义。设$p_{\bm{X}}$为真实图像分布，复原规则$p(\widehat{\bm{X}}\mid\bm{Y})$诱导输出分布$p_{\widehat{\bm{X}}}$，以$\Delta$度量逐对失真，以$d$度量两种分布的差异。给定失真预算$D$，定义
\begin{equation}
 P(D)=
 \inf_{\substack{p(\widehat{\bm{X}}\mid\bm{Y})\\
  \mathbb{E}[\Delta(\bm{X},\widehat{\bm{X}})]\leq D}}
 d(p_{\bm{X}},p_{\widehat{\bm{X}}})\text{。}
 \label{eq:vision-perception-distortion}
\end{equation}
放宽$D$扩大可选复原规则集合，因此$P(D)$不增；若$d$对第二个分布具有适当凸性，还可得到凸的最优边界。这个框架中的感知差异比较分布，不能直接替换成任意一个单图美观分数\scite{vision-perception-distortion2018}

一个完全没有信息的观测也能说明差别。假设原像素以相同概率为0或1，平方误差最优输出为0.5，平均失真0.25，但输出分布只有一个中间值。若独立地以相同概率输出0或1，边缘分布与真实像素相同，平均平方误差却为0.5。后一规则没有因此猜对每个原像素。实际方法远离最优边界时，改进模型仍可能同时降低多项误差；上述取舍不表示所有评价指标必然互相恶化。

\subsection{下游用途与信息保留}
\label{sec:vision-restoration-downstream}

招牌任务可以固定一套文字读取器，分别输入原照片、处理结果和可信参考，再与人工核实的字符串比较。若原图识别成“18路”、处理后成“13路”、参考为“18路”，锐化后的感知改善没有帮助这次读取。还应检查处理图像本身是否改出了“3”，或字符仍为“8”但读取器产生误判，以区分图像处理与下游系统的误差。

对于测量，可固定边缘定位与标定步骤，比较长度、角点和深度；对于识别，可比较相同对象的类别及拒识结果。阈值需要在验证集上确定，不能反复查看测试标签后为某个处理器调优。若联合训练了处理器与下游模型，还应把这看作一个新系统，与使用固定下游模型的对照分别报告。

当真实参考不可得时，可以保留原始RAW、多次曝光或其他视角作为额外证据，并记录生成补全区域。补出的皮肤、纹理和字符是否可接受，由用途决定；需要记录保真的任务应能追溯处理前后对应关系，不能以模型确信度替代丢失观测。

\subsection{质量与资源的取舍}
\label{sec:vision-restoration-resources}

复原比较应固定输入分辨率、退化、硬件、采样步数和候选数量。前馈网络与扩散流程的成本不只相差参数量，后者还需要多次迭代和可能的引导反向计算。图像压缩则同时统计实际文件码率、编码时间、解码时间与内存，不能只比较训练的估计码率。

表\ref{tab:vision-restoration-choice}把三类目的与所需证据对应起来。一个应用可以同时占据多行，例如先从RAW复原，再调整显示并压缩发送。最终判断仍需回到整个流程的输出，前段留下的错误可能被后段放大。

\begin{table}[htbp]
 \centering
 \small
 \begin{tabularx}{\linewidth}{@{}p{18mm}XX@{}}
  \toprule
  主要目的 & 需要保留的证据 & 共同记录的成本与条件\\
  \midrule
  内容复原 & 对齐参考、退化回投影、文字或测量结果 & 输入范围、模型版本、迭代与候选预算\\
  呈现增强 & 原始图像、亮部裁剪、颜色与局部层次变化 & 颜色空间、曲线、增益和局部参数\\
  存储传输 & 独立解码结果、实际文件、用途误差 & 辅助比特、元数据、编码及解码资源\\
  \bottomrule
 \end{tabularx}
 \caption{按处理目的选择证据。表中各行可以组合；每项判断都同时保留输入条件、处理结果与成本，避免只用一个质量分数作选择。}
 \label{tab:vision-restoration-choice}
\end{table}

现实退化覆盖不足、生成细节缺少可核验依据、低码率下关键文字丢失，分别指向数据建模、观测约束和用途相关编码的问题。研究这些问题时，应明确减少了哪种错误、付出了多少成本，以及结论依赖哪类输入。平均感知质量上升若伴随关键字段错误增加，就需要改变目标或应用范围。

\section{本章小结}
\label{sec:vision-restoration-summary}

让夜间招牌变得可用，需要区分看清、看着自然和记录正确。残差监督与显式反卷积利用不同依据估计内容，生成先验为不可辨识区域提供可能细节，RAW映射和色调调整又使用不同阶段的信息。观测丢失的细节不会因为输出更锐利而自动获得真实性。

保存结果还要经过量化、概率分配和真实熵编码，辅助信息必须同时进入码率与解码依赖。选择流程时，应把参考失真、感知变化、任务结果和资源放在同一组条件下比较。下一章允许根据用户要求创建和修改内容，届时仍需明确哪些变化是目标、哪些内容必须保持。
